\documentclass[11pt,a4paper]{article}

% ===== XeLaTeX + Korean =====
\usepackage{kotex}
\usepackage{fontspec}
\setmainfont{Times New Roman}
\setmainhangulfont{AppleMyungjo}
\setsanshangulfont{Apple SD Gothic Neo}

% ===== Layout =====
\usepackage[a4paper,top=18mm,bottom=18mm,left=20mm,right=20mm]{geometry}
\usepackage{fancyhdr}
\setlength{\headheight}{24pt}
\usepackage{enumitem}
\usepackage{mdframed}
\usepackage{array}

\setlength{\parindent}{1em}
\linespread{1.18}

% ===== Header / Footer =====
\pagestyle{fancy}
\fancyhf{}
\renewcommand{\headrulewidth}{0.6pt}
\fancyhead[L]{\sffamily\bfseries 2026 고1 영어 변형 — 지문 38}
\fancyhead[R]{\sffamily 제스처와 기억}
\renewcommand{\footrulewidth}{0pt}
\fancyfoot[C]{\framebox[16mm][c]{\thepage}}

% ===== helpers =====
\newcommand{\qno}[1]{\par\vspace{8pt}\noindent{\sffamily\bfseries #1.}\ }
\newcommand{\instr}[1]{{\sffamily #1}\par\vspace{3pt}}
\newlist{choices}{enumerate}{1}
\setlist[choices]{label=,leftmargin=1.5em,itemsep=2pt,topsep=3pt,parsep=0pt}

\newmdenv[linewidth=0.6pt,roundcorner=3pt,innertopmargin=7pt,innerbottommargin=7pt,
          innerleftmargin=9pt,innerrightmargin=9pt,skipabove=5pt,skipbelow=5pt]{pbox}
\newcommand{\nemo}[1]{\fbox{\,#1\,}}

\begin{document}

\begin{center}
{\fontsize{15}{17}\selectfont\sffamily\bfseries [지문 38] 다음 글을 읽고 물음에 답하시오.}
\end{center}
\vspace{4pt}

% ===== 삽입 문장 박스 =====
\begin{pbox}
\textbf{주어진 문장:} This suggests that the act of gesturing itself, not just following instructions to gesture, plays a role in strengthening memory.
\end{pbox}

\vspace{4pt}

% ===== 공통 지문 (문장 삽입 번호 + 빈칸) =====
\begin{pbox}
\hspace{1em}If you gesture while describing an event, do you think you'll be able to recall that event better than if you don't gesture? ( ① ) We had adults watch videos of toy objects, animals, and people performing various, sometimes odd, actions: a chicken sliding to a policeman, a jogger bending down to touch his toes, a fence swinging shut on its own. ( ② ) We then tested the adults' memories of these events immediately after the descriptions and three weeks later. ( ③ ) To see if we'd get the same effects if people gestured on their own, we did the study again, but this time we gave the adults no instructions about gesture. ( ④ ) We found the same patterns: people remembered items on which they had gestured \rule{3.5cm}{0.4pt} better than items on which they had not gestured. ( ⑤ ) Producing gesture along with speech makes the information encoded in that speech memorable.

\vspace{4pt}
\noindent\footnotesize{* encode: 부호화하다}
\end{pbox}

% ===================== Q1 문장 삽입 =====================
\qno{1}\instr{글의 흐름으로 보아, 주어진 문장이 들어가기에 가장 적절한 곳을 고르시오.}
\begin{choices}
\item ① \quad ② \quad ③ \quad ④ \quad ⑤
\end{choices}

% ===================== Q2 빈칸추론 =====================
\qno{2}\instr{다음 빈칸에 들어갈 말로 가장 적절한 것은?}
\begin{choices}
\item ① in response to direct instruction
\item ② of their own accord
\item ③ at a rapid pace
\item ④ in a repetitive manner
\item ⑤ with great exaggeration
\end{choices}

% ===================== Q3 서술형 요지 =====================
\qno{3}\instr{[서술형] 이 글의 요지를 한 문장(40자 이내)으로 우리말로 쓰시오.}
\par\vspace{6pt}\noindent
$\Rightarrow$ \rule{10cm}{0.4pt}

% ===================== Q4 =====================
\qno{4}\instr{윗글의 내용과 일치하지 \underline{않는} 것은?}
\begin{choices}
\item ① Adults in the study watched videos of various unusual actions.
\item ② Memory was tested both immediately and several weeks later.
\item ③ A second study examined people who gestured without explicit instruction.
\item ④ The study found no memory difference between gestured and non-gestured items.
\item ⑤ Producing gestures with speech can make spoken information more memorable.
\end{choices}

\end{document}
